\section{Začínáme}

\begin{frame}{Co je kde}
  \framesubtitle{Obsah se mění často, motiv skoro vůbec}
  \begin{columns}[T,totalwidth=\textwidth]
    \begin{column}{.48\textwidth}
      \textbf{Pro každou prezentaci}
      \begin{itemize}
        \item \texttt{presentation-config.tex}\\ název, autor, datum a titulka
        \item \texttt{slides/*.tex}\\ jednotlivé kapitoly a slidy
        \item \texttt{assets/}\\ loga, fotografie a diagramy
      \end{itemize}
    \end{column}
    \begin{column}{.48\textwidth}
      \textbf{Sdílený vzhled}
      \begin{itemize}
        \item \texttt{beamerthemeVSAK.sty}\\ titulka, hlavička a patička
        \item \texttt{vsak-colors.sty}\\ jediný zdroj barev
        \item \texttt{vsak-components.sty}\\ boxy a pomocné prvky
      \end{itemize}
    \end{column}
  \end{columns}
  \begin{Tip}
    Běžně stačí upravovat konfiguraci, soubory ve \texttt{slides/} a vlastní obrázky.
  \end{Tip}
\end{frame}

\begin{frame}{Jak vznikne prezentace}
  \framesubtitle{Čtyři kroky od kopie šablony k hotovému PDF}
  \begin{enumerate}
    \item Uprav metadata v \texttt{presentation-config.tex}.
    \item Rozděl obsah do tematických souborů v \texttt{slides/}.
    \item Vlož obrázky do \texttt{assets/} a použij relativní cestu.
    \item Sestav prezentaci z hlavního souboru \texttt{main.tex}.
  \end{enumerate}
  \vspace{4mm}
  \begin{Rule}
    Jeden soubor se slidy = jedna kapitola. Novou kapitolu připoj jediným příkazem \texttt{\textbackslash input} v \texttt{main.tex}.
  \end{Rule}
\end{frame}
